\chapter*{Group Project: Integrated eNSP Network}
\addcontentsline{toc}{chapter}{Group Project: Integrated eNSP Network}\markboth{Group Project}{}
\begin{sourcebasis}Based on \textbf{Group Project - Template DCN.ppt}.\end{sourcebasis}

\section*{What the project asks you to demonstrate}
Build a client-server scenario in eNSP and explain it through \textbf{two or three OSI/TCP-IP layers}.

\begin{tabularx}{\linewidth}{>{\bfseries}p{0.22\linewidth}X}
\toprule
Layer/theme & Project task from the supplied template\\\midrule
Application & Configure a client-server service in eNSP.\\
Transport & Demonstrate TCP connection setup, rate limiting, packet loss, and retransmission.\\
Network & Build 2--3 routers with multiple hosts, IPv4 addressing, static routing, and packet forwarding.\\
Data Link (optional) & Explain how frames cross the selected physical or logical link.\\\bottomrule
\end{tabularx}

\section*{A build sequence that reuses the lab workflow}
\begin{netsteps}
\item Draw the topology and write the addressing plan before entering any CLI commands.
\item Configure host/router interfaces and prove every directly connected link.
\item Add static routes and verify each remote destination in the routing table.
\item Prove end-to-end forwarding before adding the client-server service.
\item Add the chosen application/TCP behavior and collect packet or command evidence.
\item If using Data Link, show its effect on the end-to-end flow.
\item Label each readable screenshot and tie it to a technical claim.
\end{netsteps}

\section*{Recommended evidence matrix}
\begin{tabularx}{\linewidth}{>{\bfseries}p{0.26\linewidth}X}
\toprule
Claim in presentation/report & Strong evidence\\\midrule
``The topology is addressed correctly'' & Addressing table plus \cmd{display ip interface brief}/host IP configuration.\\
``Routers can forward between networks'' & Routing-table entries and successful end-to-end ping.\\
``TCP provides reliable transfer'' & Capture of connection setup and retransmission after loss.\\
``Rate limiting affects the flow'' & Before/after measurements tied to the configured limit.\\
``The design uses a specific layer/protocol'' & Configuration excerpt plus a packet/state observation that proves the protocol is active.\\\bottomrule
\end{tabularx}

\section*{Deliverables}
The project has three outputs:
\begin{itemize}
\item \textbf{eNSP demo:} a working network explained while it runs.
\item \textbf{Presentation:} design, key configuration, protocol choices, and evidence.
\item \textbf{Report:} a record of the implementation. The source states Week 12 submission through uLearn and to the lab instructor.
\end{itemize}

\section*{Evaluation criteria}
The rubric weights are:
\begin{center}\begin{tabular}{lr}\toprule
Design \& Implementation & 40\%\\
Functionality \& Demo & 30\%\\
Presentation \& Explanation & 20\%\\
Report Quality & 10\%\\\bottomrule
\end{tabular}\end{center}
Bands are Excellent (90--100), Good (75--89), Satisfactory (60--74), and Needs Improvement (<60).

\sourceimage[0.74\linewidth]{Evaluation-criteria table from the supplied project template}{project-evaluation.png}

\begin{keyidea}
Pair each claim with evidence: interface state, a route, a packet trace, an application result, or a connectivity test.
\end{keyidea}

\begin{labdeliverable}
Submit a topology, addressing table, configuration excerpts, routing and TCP evidence, concise layer explanations, and brief troubleshooting/limitations notes.
\end{labdeliverable}
